\section{Findings}
\label{sec:findings}

Each finding below is stated in bold, followed by the evidence already
presented and the bound on what it establishes. Nothing appears here that
cannot be read off a table or figure earlier in the report, and nothing is
stated more strongly than the sample behind it allows.

\paragraph{\textbf{On recordings it never saw, the diffusion family beats the
action-chunking family --- by far less than the obvious comparison suggests.}}
Over the horizon the two share, the margin is under a tenth; over their own
horizons it appears to be roughly a third. The direction survives the
correction and the magnitude does not. \emph{Bound:} the two families differ in
architecture and in planning horizon at once, and only the second has been
controlled for here.

\paragraph{\textbf{The two families differ more in how much they degrade than
in how well they perform.}} On familiar recordings the action-chunking family
is the more accurate of the two; on unfamiliar ones it is the less accurate. Its
error grows roughly fourfold between the two halves against under twofold for
the diffusion family. \emph{Bound:} this is the largest and most robust effect
in the report, being a ratio within each policy and immune to the horizon
problem --- but the held-out recordings are the last of the session, so part of
it may be drift rather than overfitting.

\paragraph{\textbf{Cropping the tactile rim buys nothing that this experiment
can detect.}} On held-out recordings both families move by about one per cent
or less, in opposite directions, on a sample too small to resolve either.
\emph{Bound:} a null result for \emph{this} crop, which removes the rim and
stretches the remainder together; it is not evidence that the rim carries no
information.

\paragraph{\textbf{The crop's effect changes sign with horizon.}} For the
action-chunking family, cropping is better over the first thirty-two steps of
the chunk and worse over the full hundred. \emph{Bound:} this is visible only
because the error is reported as a curve; both numbers come from the same pair
of runs, which is the point.

\paragraph{\textbf{Neither family over-reads the gel rim, and cropping moves
attention toward it.}} Gradient mass in the outer band sits at two thirds to
four fifths of what an indifferent model would place there, at every band width
tested, and rises for both families after cropping. This contradicts the
observation that motivated the crop. \emph{Bound:} measured on training
recordings, and the rise after cropping is partly mechanical, since removing
the outer rows moves the surviving pixels closer to the new edge.

\paragraph{\textbf{The families do not use the rig the same way, and the
tactile channels rank last for two of the three.}} The action-chunking family
leans on the overhead camera, the diffusion family on proprioception, and the
third ranks touch above proprioception --- the only one to do so.
\emph{Bound:} shares are normalised within a policy, so this is a claim about
ordering and never about magnitudes; and the third family reads half as many
tactile cameras as the others.

\paragraph{\textbf{A policy can be flat across its entire planning horizon on
recordings it trained on.}} The action-chunking family predicts the hundredth
step of a familiar recording about as accurately as the first. That is what
reproducing a memorised trajectory looks like, and it is not what planning
looks like. \emph{Bound:} a statement about behaviour on training data, not
about competence on the rig.

The two findings that bear most on what to do next are the second and the
third. The generalisation gap is large, robust, and currently confounded with
session drift; the crop --- the change actually under test --- does nothing
measurable. Together they point away from preprocessing the tactile input and
toward the amount and variety of data behind these policies.
